\documentclass{tutodoc}

\usepackage{../preamble.cfg}

\usepackage{../main/main}


\begin{document}

%\section{La syntaxe du langage proposé}
%\subsection{Les contextes pour les signes et/ou les variations}

\subsubsection{Le contexte \tdoclatexin{vars} pour les variations}

La logique d'utilisation et d'organisation des informations pour le contexte \tdoclatexin{vars} est similaire à celle du contexte \tdoclatexin{signs}\,, nous indiquons donc juste les différences.
%
\begin{enumerate}
    \item Voici les informations possibles pour le comportement d'une expression $f$ en supposant que $n$ valeurs pivots \tdoclatexin{x1}\,, \tdoclatexin{x2}\,, \dots\ , \tdoclatexin{xn} ont été données via \tdoclatexin{xvals = x1 , x2 , ... , xn}\,. De nouveau, nous posons $p = n - 1$\,.
    %
    \begin{itemize}
	    \item Un seul contexte \tdoclatexin{vars} peut ne pas nommer l'expression pour utiliser $f$ le nom par défaut, les autres devront le faire via \tdoclatexin{monexpr : ...} où les points de suspension spécifient les informations liées au comportement de l'expression (voir ce qui suit).


        \item A minima, il faut indiquer
        \tdoclatexin{v1 v2 ... vp}
        où
        \tdoclatexin{v1} , \tdoclatexin{v2} , \dots\ , \tdoclatexin{vp}
        donnent des informations sur les $p$ intervalles ouverts
        \tdoclatexin{]x1 ; x2[} , \tdoclatexin{]x2 ; x3[} , \dots\ , \tdoclatexin{]xp ; xn[} respectivement.
        Les valeurs possibles pour les \tdoclatexin{vk} sont les suivantes.
        %
        \begin{enumerate}
            \item \tdoclatexin{<} indique une expression strictement croissante sur l'intervalle concerné.

            \item \tdoclatexin{>} indique une expression strictement décroissante sur l'intervalle concerné.

            \item \tdoclatexin{=} indique une expression constante sur l'intervalle concerné.

            \item \tdoclatexin{u} indique une expression non définie sur l'intervalle concerné (comme pour les signes).
        \end{enumerate}


        \item Pour des valeurs pivots précises, on peut indiquer des images ou des limites à droite et/ou à gauche. Voici ce qui est disponible.
        %
        \begin{enumerate}
            \item L'absence d'expression est possible pour ne rien indiquer du tout.

            \item Toute expression sans lettre \tdoclatexin{u}, ni ponctuation \tdoclatexin{!} est interprétée comme une valeur image au format mathématique \LaTeX.

            \item Une valeur image contenant \tdoclatexin{u} et/ou la ponctuation \tdoclatexin{!} en tant que \tdocquote{token \LaTeX} devra être protégée par des accolades. Ceci vient du fait que dans le langage codant les variations ces deux caractères ont une signification spéciale (voir ci-dessus et l'item suivant).

            \item \tdoclatexin{!} indique que
l'expression n'est pas définie au pivot concerné.
            On peut aussi, si besoin, indiquer des limites à gauche et/ou à droite. Voici les cas possibles.
            %
            \begin{itemize}[label={\scriptsize o}]
                \item \tdoclatexin{!} utilisé seul n'indique aucune limite.


                \item  \tdoclatexin{! d} indique juste $d$ comme limite à droite.

                \textbf{Cette syntaxe est interdite pour le tout dernier pivot.}


                \item \tdoclatexin{g !} indique juste $g$ comme limite à gauche.

                \textbf{Cette syntaxe est interdite pour le tout premier pivot.}


                \item  \tdoclatexin{g ! d} indique $g$ et $d$ comme limites à gauche et à droite respectivement.

                \textbf{Cette syntaxe est interdite pour les premier et dernier pivots.}
            \end{itemize}
        \end{enumerate}


    	\item \textbf{Il est possible de renseigner une ligne partiellement.}
    \end{itemize}
\end{enumerate}


% --------------- %


%Voici un tableau permettant de mémoriser l'organisation des informations de type \tdocquote{variation} et celles de type \tdocquote{pivot}.


% --------------- %


\begin{tdocnote}
    Les espaces autour des doubles points, du point d'exclamation et des informations codées ne sont pas obligatoires.
\end{tdocnote}


% --------------- %


Voici deux codes fictifs illustrant les explications précédentes ; noter au passage que les espaces ignorés permettent d'obtenir un résultat humainement très clair.
\begin{multicols}{2}
    \tdoclatexinput<\tdoctcb{code}>{examples/dsl-l3/vars-always-defined.tex}

    \tdoclatexinput<\tdoctcb{code}>{examples/dsl-l3/vars-with-undefined.tex}
\end{multicols}

Dans le cadre de processus automatisés, il est possible de produire les horreurs suivantes qui aboutiront aux mêmes sorties que les codes correspondants ci-dessus.
\begin{multicols}{2}
    \tdoclatexinput<\tdoctcb{code}>{examples/dsl-l3/vars-always-defined-monster.tex}

    \tdoclatexinput<\tdoctcb{code}>{examples/dsl-l3/vars-with-undefined-monster.tex}
\end{multicols}

\end{document}
